\documentclass[10pt,a4paper]{amsart}
\usepackage[margin=19mm]{geometry}
\usepackage{amsmath,amssymb,mathtools}
\usepackage[scr=rsfs,cal=euler]{mathalfa}
\usepackage{xfrac}
\usepackage[unicode,hidelinks,bookmarks=false]{hyperref}
\newcommand{\ie}{i.e.\ }
\newcommand{\cf}{cf.\ }
\newcommand{\resp}{respectively\ }
\newcommand{\Subsection}{\subsection}
\providecommand{\nobreakdash}{\nobreak\hbox{-}}
\setlength{\emergencystretch}{2em}
\multlinegap0pt
\newtheorem{theorem}{Theorem}
\newcommand{\Z}{{\mathbb Z}}
\newcommand{\s}{{\mathbb S}}
\title{Small symplectic $4$-manifolds via contact gluing and some applications}
\author{Weimin Chen}
\date{Source: arXiv:2503.05932v6 (CC BY 4.0)}
\begin{document}
\maketitle

\noindent\emph{Source and licence.} Small symplectic $4$-manifolds via contact gluing and some applications. Version arXiv:2503.05932v6 (CC BY 4.0), distributed by arXiv under the Creative Commons Attribution 4.0 International licence, http://creativecommons.org/licenses/by/4.0/, as recorded in the arXiv licence metadata for this exact version and shown on the abstract page https://arxiv.org/abs/2503.05932v6. Full licence notice: \texttt{licenses/arxiv-2503.05932-CC-BY-4.0.txt}.\par\smallskip

\noindent\textit{Editorial scope.} Counted unit: the article's theorem theorem (source0.tex lines 1245--1264). The article's complete original proof of it is delivered with it (source0.tex lines 1266--1356, one \texttt{proof} environment of 91 lines). No mathematics is written, repaired or re-derived: the statement and the proof are the article's own lines, in the article's own notation, and no retained line is substituted or re-flowed.  No further source-local context is needed: every cross-reference of every retained line resolves inside the two delivered spans.  Nothing was omitted from any retained span, so the package carries no omission record.  No retained line contains a citation, so the wrapper carries no bibliography and no bibliography entry.  No retained line contains a figure environment, an image-inclusion command or drawing code, so no image is delivered and the package declares no asset dependency.  Editorial formatting only, in addition: an AMS \texttt{amsart} preamble that restates the article's own declarations and macros used by the retained lines, namely the environments \texttt{theorem} on separate counters, together with the standard packages those lines need.  The retained environments print in this document as Theorem 1 in order of appearance, whereas the article prints the same items with the numbering of its own full document.  The complete native source of the article as distributed in the arXiv e-print for this exact version is bundled unchanged as \texttt{sources/174429/source0.tex}.
\begin{theorem}
Assume $M$ is a Seifert manifold, i.e., the $\s^1$-action induced by the rational open book 
$(B,\pi)$ is fixed-point free, and as a Seifert manifold, $M$ is given by 
$$
M=M(g_0; \{(\alpha_i,\beta_i)\}, \{(\alpha_j,\beta_j)\}),
$$
where $0<\beta_j<\alpha_j$ for each $j$, and $\{(\alpha_i,\beta_i)\}$ are the Seifert invariants at
the binding components $B_i$, subject to the following convention: if $\alpha_i<0$ for some $i$, 
then $B_i$ has the opposite orientation of the fiber.  Let $n$ be the order of the monodromy of 
$(B,\pi)$ and let $p_i$ be the multiplicity at the binding component $B_i$. 
Then $M^\prime$ is given by
$$
M^\prime=M(g_0; \{(\bar{\alpha}_i,\bar{\beta}_i)\}, \{(\alpha_j,-\beta_j)\}),
$$
where $\bar{\alpha}_i=\frac{n}{p_i}-F_i\alpha_i$ and 
$\bar{\beta}_i=F_i\beta_i+\frac{1}{\alpha_i}-\frac{n\beta_i}{p_i\alpha_i}$, subject to the convention that
if $\bar{\alpha}_i<0$, then the corresponding binding component $B_i^\prime$ of the rational open 
book $(B^\prime,\pi^\prime)$ has the opposite orientation of the fiber. Finally, $\bar{\alpha}_i=0$
is also allowed here.  
\end{theorem}

\begin{proof}
The proof is simply an application of Lemma 3.2.

To this end, we begin by recalling the standard description of $(B,\pi)$ used in the proof of Lemma 3.2. First, the periodic monodromy map $\phi: \Sigma\rightarrow \Sigma$ determines a
set of integers $\{(n_j,c_j)\}$, $\{(n,c_i)\}$, where $0<c_j<n_j$, $gcd(n_j,c_j)=1$, 
$0<c_i<n$, $gcd(n,c_i)=1$, and moreover, for each $i$, there are integers $k_i,l_i$, uniquely determined by $c_i$ by the conditions $0<k_i<n$, $nl_i-k_ic_i=1$. Finally, 
$$b:=\sum_j\frac{c_j}{n_j}+\sum_i\frac{c_i}{n}\in\Z.$$

To describe the gluing of the regular neighborhood $N_i$ of $B_i$ to the mapping torus 
$\Sigma\times [0,1]/ (x,1)\sim (\phi(x),0)$, we fix a canonical, positively oriented basis 
$(\gamma_i,\tau_i)$ of $H_1(T_i)$, where $T_i:=(\partial \Sigma)_i \times [0,1]/ (x,1)\sim (\phi(x),0)$,
and $\gamma_i:=(\partial \Sigma)_i \times \{0\}$. With this understood, the meridian $\mu_i$ of
$\partial N_i=-T_i$ is given by $\mu_i=-p_i^\prime\gamma_i +p_i\tau_i$ for a pair of integers 
$(p_i,p_i^\prime)$, where $p_i>0$ is the multiplicity at $B_i$.

With $(p_i,p_i^\prime)$ understood, we note that $\alpha_i=k_ip_i+np_i^\prime$. On the other hand,
we set 
$$
\beta_i^\prime:=l_ip_i+c_ip_i^\prime.
$$
Then by Lemma 3.2, $M=M(g_0; (1,-b), \{(\alpha_j,\beta_j)\}, \{(\alpha_i,\beta_i^\prime)\})$
where $(\alpha_j,\beta_j)=(n_j,c_j)$ for each $j$.  Observe that $\sum_i\frac{\beta_i}{\alpha_i}=\sum_i\frac{\beta_i^\prime}{\alpha_i}-b$.

With the preceding understood, we shall apply Lemma 3.2 to the rational open book 
$(B^\prime,\pi^\prime)$ to obtain the generalized Seifert fibration structure on $M^\prime$. To this end,
we identify $\pi^\prime: M^\prime\setminus B^\prime \rightarrow \s^1$ with $\pi: M\setminus B \rightarrow \s^1$ but with the orientation of the pages $\Sigma$ reversed. We shall describe the rational open book $(B^\prime,\pi^\prime)$ in the same fashion, using the bar-version of the notations. Since the page orientation is reversed, the set of integers $\{(n_j, \bar{c}_j)\}$, $\{(n,\bar{c}_i)\}$ determined by the monodromy map are given by $\bar{c}_j=n_j-c_j$, $\bar{c}_i=n-c_i$. 
Consequently, $\bar{k}_i=n-k_i$, $\bar{l}_i=n+l_i-k_i-c_i$, and 
$$
\bar{b}=\sum_j\frac{\bar{c}_j}{n_j}+\sum_i\frac{\bar{c}_i}{n}=x+y-b,
$$
where $x,y$ are the number of indices of $j$ and $i$ respectively. Finally, we note that the corresponding basis $(\bar{\gamma}_i,\bar{\tau}_i)$ is related to $(\gamma_i,\tau_i)$ by the 
equations $(\bar{\gamma}_i,\bar{\tau}_i)=(-\gamma_i, \tau_i+\gamma_i)$. 

In order to obtain the generalized Seifert invariants of $M^\prime$, it remains to determine the
corresponding pair of integers $(\bar{p}_i,\bar{p}_i^\prime)$. To this end, 
we fix a longitude $\lambda_i$ along $B_i$, and write $\gamma_i=p_i\lambda_i+q_i\mu_i$, 
$\tau_i=p_i^\prime\lambda_i+q^\prime_i\mu_i$ for some $q_i,q_i^\prime$ such that 
$p_iq_i^\prime-p_i^\prime q_i=1$ (note that $\lambda_i\cdot\mu_i=1$). 
Suppose with respect to $(\lambda_i,\mu_i)$, the framing of the $2$-handle attached along $B_i$
is $f_i$. Then the framing $f_i$ and the framing relative to the page framing $F_i$ are related by
the equation $f_i=F_i+\frac{q_i}{p_i}$. On the other hand, as a boundary component of $Z$, 
$-M^\prime$ is obtained by removing the neighborhood $N_i$ of $B_i$ in $\{1\}\times M$ and then
gluing back a $D^2\times \s^1$ along $\partial N_i$ such that 
$$
\mu_i^\prime=-\lambda_i-f_i\mu_i,\;\; \lambda_i^\prime=\mu_i,
$$
where $(\mu_i^\prime,\lambda_i^\prime)$ is the meridian-longitude pair of the solid torus 
$D^2\times \s^1$. 

With the preceding understood, we observe that the reversing of the orientation of the pages $\Sigma$ 
results a change of the orientations of $(\mu_i^\prime,\lambda_i^\prime)$ to the opposite. Consequently, if we let $\bar{\mu}_i=-\mu_i^\prime$ and $\bar{\lambda}_i=-\lambda_i^\prime$, 
then we shall determine the pair $(\bar{p}_i,\bar{p}_i^\prime)$ through the equation
$\bar{\mu}_i=-\bar{p}_i ^\prime\bar{\gamma}_i+\bar{p}_i \bar{\tau}_i$. On the other hand, 
the reversing of the orientation of the pages $\Sigma$ also results a change of the orientation of
$M$ to the opposite, so that the intersection numbers $\gamma_i\cdot\tau_i$ and 
$\lambda_i\cdot \mu_i$ are changed by a sign, i.e., we now have 
$\gamma_i\cdot\tau_i=\lambda_i\cdot \mu_i=-1$. (Note that $\bar{\gamma}_i\cdot \bar{\tau}_i=1$
is unchanged.) Finally, note that $\mu_i^\prime=-\lambda_i-f_i\mu_i$ implies that 
$\bar{\mu}_i=\lambda_i+f_i\mu_i$. With this understood, and appealing to the relation
$(\bar{\gamma}_i,\bar{\tau}_i)=(-\gamma_i, \tau_i+\gamma_i)$, we compute $\bar{p}_i$, 
$\bar{p}_i^\prime$ through intersection numbers:
$$
\bar{p}_i=\bar{\gamma}_i\cdot \bar{\mu}_i=p_if_i-q_i, \;\;\; 
\bar{p}_i^\prime=-\bar{\mu}_i\cdot \bar{\tau}_i=(q_i+q_i^\prime)-f_i(p_i+p_i^\prime).
$$
Note that $\bar{p}_i=p_i F_i$. In particular, $F_i>0$ implies $\bar{p}_i>0$.

Set $\bar{\alpha}_i:=\bar{k}_i\bar{p}_i+n\bar{p}_i^\prime$, 
$\bar{\beta}_i^\prime:=\bar{l}_i\bar{p}_i+\bar{c}_i\bar{p}_i^\prime$. Then with the relations
$$
F_i=f_i-\frac{q_i}{p_i},\;\; p_iq_i^\prime-p_i^\prime q_i=1, \;\; \alpha_i=k_ip_i+np_i^\prime, \;\;\beta_i^\prime=l_ip_i+c_ip_i^\prime, 
$$
a straightforward calculation shows that
$$
\bar{\alpha}_i=\frac{n}{p_i}-F_i\alpha_i, \;\;\;
\bar{\beta}_i^\prime=F_i\beta_i^\prime+\frac{1}{\alpha_i}
-\frac{n\beta_i^\prime}{p_i\alpha_i}+\bar{\alpha}_i.
$$
Now if we write $\beta_i^\prime=\beta_i+b_i\alpha_i$, then
$$
\bar{\beta}_i^\prime=F_i\beta_i+\frac{1}{\alpha_i}
-\frac{n\beta_i}{p_i\alpha_i}+ (1-b_i)\bar{\alpha}_i.
$$
With $\bar{b}=x+y-b$, $\sum_ib_i=b$, and $(n_j,c_j)=(\alpha_j,\beta_j)$, it follows from Lemma 3.2 that
$$
M^\prime=M(g_0; (1,-\bar{b}), \{(n_j,\bar{c}_j)\}, \{(\bar{\alpha}_i,\bar{\beta}_i^\prime)\})=
M(g_0; \{(\bar{\alpha}_i,\bar{\beta}_i)\}, \{(\alpha_j,-\beta_j)\}),
$$
where $\bar{\beta}_i=F_i\beta_i+\frac{1}{\alpha_i}-\frac{n\beta_i}{p_i\alpha_i}$. This finishes the proof of Theorem 4.4.

\end{proof}

\end{document}
